\subsection{Cyclic level structures and new filtered height-two counterexamples}

\begin{lemma}[A tame equivariant arc with a rational cyclic level structure]
\label{lem:B-RD62-arc}
\status{Proved; independent audits passed.}
Let $p\ge5$, let $n\ge2$, and let $e$ be a positive multiple of $p-1$
such that
\[
                    p^{2n-3}<e<p^{2n-2}.                 \tag{B-RD62.1}
\]
There is a complete discretely valued field $K$ of mixed characteristic
$(0,p)$, residue field $\overline{\mathbf F}_p$, and absolute ramification
$e$, containing $\mu_p$, and a height-two BT group $G/K$ with a
$K$-rational point of exact order $p^n$.  It has a good one-dimensional
formal model over a cyclic tame extension $L/K$ of degree $p^2-1$.
On its covariant special isocrystal, a generator of this descent group
has eigenvalues $\zeta^{-1},\zeta^{-p}$, where $\zeta$ is a primitive
$(p^2-1)$st root of unity.
\end{lemma}
\begin{proof}
Put $k=\overline{\mathbf F}_p$, $W=W(k)$, $q=p^2$, $N=q-1$, and
\[
             d=Ne,\qquad b=d-(p-1)q^{n-1}.              \tag{B-RD62.2}
\]
The inequalities in (B-RD62.1) imply
\[
 \begin{gathered}
 b>(q-p)q^{n-2}>0,\qquad q^{n-1}<d<Nq^{n-1},\\
 b\equiv-(p-1)\pmod N.
 \end{gathered}                                         \tag{B-RD62.3}
\]
For the first inequality, subtracting its right side from $b$ gives
$N(e-p^{2n-3})$.  The last congruence uses $q\equiv1\pmod N$.

Use the integral $p$-typical group $F_u$ over $W[[u]]$ obtained by
specializing $v_1=u$, $v_2=1$, and $v_j=0$ for $j>2$ in
\cite[Proposition~5.7, equations (5.3)--(5.5), and (12.3)--(12.5)]
{HopkinsGross1994Bundles}.
Its logarithm is
\[
 f_u(T)=\sum_{j\ge0}l_j(u)T^{p^j},\qquad l_0=1,
 \qquad p l_j=l_{j-1}u^{p^{j-1}}+l_{j-2},               \tag{B-RD62.4}
\]
where $l_{-1}=0$.  This is precisely the specialized recursion (5.5),
not a claim that the multiplication series itself is a polynomial.
The same source gives
\[
 \begin{gathered}
 [p]_u(T)=pT+(1-p^{p-1})uT^p+\sum_{j>p}c_j(u)T^j,\\
 c_j(u)\in(p,u)W[[u]]\quad(p<j<q),\\
 c_q(u)\equiv1\pmod{(p,u)}.
 \end{gathered}                                         \tag{B-RD62.5}
\]
The coefficients in degrees $2,\ldots,p-1$ vanish.  The special group
$F_0\otimes k$ is the height-two Honda group.

Define integral power series in $W[[X]]$ by
\[
 \begin{split}
 Y_0(X)&=X,\\
 Y_{j+1}(X)&=[p]_{X^b}\bigl(Y_j(X)\bigr),\\
 \Phi_u(T)&=[p]_u(T)/T,\\
 P(X)&=\Phi_{X^b}\bigl(Y_{n-1}(X)\bigr).
 \end{split}                                           \tag{B-RD62.6}
\]
All compositions converge in the indicated formal power-series ring.
Modulo $p$, induction using (B-RD62.3)--(B-RD62.5) gives
\[
                  Y_j(X)=X^{q^j}+O(X^{q^j+1})
                  \quad(0\le j\le n-1).                \tag{B-RD62.7}
\]
Indeed, at an input of order $q^{j-1}$ the $q$th-degree term is
strictly earlier than the $p$th-degree term because
$b>(q-p)q^{j-1}$; intermediate coefficients are multiples of $u$
modulo $p$, and all terms of degree greater than $q$ are later.
In the last series $P$, its $uY_{n-1}^{p-1}$ term has order
$b+(p-1)q^{n-1}=d$.  Every intermediate term is later, and its
$Y_{n-1}^{q-1}$ term has order $Nq^{n-1}>d$.  Consequently
\[
                    P(0)=p,\qquad
                    P(X)\equiv X^d+O(X^{d+1})\pmod p.  \tag{B-RD62.8}
\]
Weierstrass preparation therefore identifies
$R=W[[X]]/(P)$ with the ring defined by a degree-$d$ Eisenstein
polynomial.  It is a complete DVR, the image $x$ of $X$ is a
uniformizer, and $v_R(p)=d$.

The group $F_{x^b}$ has good height-two reduction over $R$.
The congruences (B-RD62.7) give $v_R(Y_j(x))=q^j$: any term whose
coefficient is divisible by $p$ has valuation at least $d>q^{n-1}$,
and the other terms are strictly later.  In particular
$Y_{n-1}(x)\ne0$.  Since $P(x)=0$,
\[
 [p]^n(x)=Y_{n-1}(x)\Phi_{x^b}(Y_{n-1}(x))=0.
\]
Thus $x$ is a point of exact order $p^n$ of this good formal group.

It remains to descend both the group and this particular point.
Choose $\zeta\in W^\times$ of order $N$ and put $a=\zeta^{-1}$.
The logarithm recursion proves, coefficient by coefficient,
\[
 l_j(a^{p-1}u)=a^{p^j-1}l_j(u),\qquad
 f_{a^{p-1}u}(T)=a^{-1}f_u(aT).                         \tag{B-RD62.9}
\]
For the $l_{j-2}$ term the exponents agree modulo $N$ because
$N\mid p^j-p^{j-2}$.  Hence multiplication of the coordinate by $a$
is an integral isomorphism $F_{a^{p-1}u}\to F_u$ and
\[
 [p]_{a^{p-1}u}(a^{-1}T)=a^{-1}[p]_u(T),\qquad
 \Phi_{a^{p-1}u}(a^{-1}T)=\Phi_u(T).                    \tag{B-RD62.10}
\]
The substitution $X\mapsto\zeta X$ sends $X^b$ to $a^{p-1}X^b$.
Induction gives $Y_j(\zeta X)=a^{-1}Y_j(X)$, and thus the exact
identity $P(\zeta X)=P(X)$.  It defines an automorphism $\sigma$ of
$R$ of order $N$, fixing $W$ and sending $x$ to $\zeta x$.

Let $L=\operatorname{Frac}(R)$ and $K=L^{\langle\sigma\rangle}$.
The extension $L/K$ is cyclic tame of degree $N$, with unchanged
residue field.  It follows that $v_K(p)=d/N=e$.
The isomorphism $\sigma^*F_{x^b}\to F_{x^b}$ with coordinate $aT$
is a descent datum: its $N$fold composite is the identity.
Moreover it fixes the marked point, since $a\sigma(x)=x$.
Effective Galois descent at every generic finite level gives the
BT group $G/K$ and its rational point.  Only the generic group is
descended; no good integral model over $\mathcal O_K$ is asserted.

On the Honda special fibre the descent is $T\mapsto aT$.
Its covariant isocrystal eigenvalues are $a,a^p$ in the conventions
of Proposition~\ref{prop:B-RD4-dictionary}; equivalently these are
the two embeddings of the unramified quadratic Teichmuller unit.
Finally $K$ contains $\mu_p$.  To see this without a polarization,
write $-p=\pi_K^e u$ with $u$ a unit.  Since $k$ is algebraically
closed, Hensel's lemma gives a $(p-1)$st root of $u$.  As $p-1\mid e$,
the field $K$ contains a $(p-1)$st root of $-p$.  Over
$\operatorname{Frac}W(k)$ this generates the unique tame extension
of degree $p-1$, which is the cyclotomic extension generated by
$\mu_p$.  No assertion that $\mu_{p^n}\subset K$ is used.
\end{proof}

\begin{lemma}[Lifting a finite rank-one extension]
\label{lem:B-RD62-extension}
\status{Proved; independent audits passed.}
Let $K$ have characteristic zero, and let
$\psi:G_K\to\mathbf Z_p^\times$ be a finite character of order
prime to $p$.  Every extension of free $\mathbf Z/p^n$-modules with
continuous Galois action
\[
 0\longrightarrow\mu_{p^n}\longrightarrow V_n
       \longrightarrow(\mathbf Z/p^n)(\psi^{-1})
       \longrightarrow0                               \tag{B-RD62.11}
\]
is the $p^n$-torsion of a height-two BT group $H/K$ admitting a
filtration with graded groups $\mu_{p^\infty}$ and
$(\mathbf Q_p/\mathbf Z_p)(\psi^{-1})$.
\end{lemma}
\begin{proof}
An underlying module splitting identifies the extension class with
$H^1(K,T/p^nT)$ for $T=\mathbf Z_p(1)\otimes\psi$.
Let $K'/K$ kill $\psi$, with group $\Delta$ of order prime to $p$.
Restriction identifies this cohomology with the $\Delta$-invariants
of $H^1(K',\mu_{p^n})\otimes\psi$, by averaging and
inflation--restriction.  The Kummer descriptions
\[
 H^1(K',\mu_{p^m})=K'^\times/(K'^\times)^{p^m}
\]
show that the transition map from level $m+1$ to level $m$ is
surjective: it sends the class of an element to the class of the
same element.  Averaging a lift by $|\Delta|^{-1}\sum_{\delta\in\Delta}\delta$
proves surjectivity also on the invariant parts.  Thus the class
at level $n$ lifts through all higher levels.

Choose continuous cocycles representing these compatible classes.
If a chosen cocycle at the next level reduces only up to a
coboundary, lift its correcting zero-cochain and subtract that
coboundary.  The cocycles can therefore be made exactly compatible.
Their inverse limit is a continuous $T$-valued cocycle and defines
an extension of free $\mathbf Z_p$-modules
\[
 0\longrightarrow\mathbf Z_p(1)\longrightarrow T_H
       \longrightarrow\mathbf Z_p(\psi^{-1})
       \longrightarrow0.
\]
The corresponding compatible finite etale group schemes over $K$
give the required BT group.  Its two rank-one graded groups acquire
good reduction after killing $\psi$, so it satisfies exactly the
monodromy condition used in B.4.  The entire extension is not claimed
to acquire good reduction.
\end{proof}

\begin{theorem}[Character ambiguity beyond minimal cyclotomic ramification]
\label{thm:B-RD62}
\status{Proved; independent audits passed.}
For every $p\ge5$, $n\ge2$, and multiple $e$ of $p-1$ satisfying
(B-RD62.1), there are height-two BT groups $M,H$ over a field $K$
with the properties in Lemma~\ref{lem:B-RD62-arc}, such that
\[
                            M[p^n]\simeq H[p^n],        \tag{B-RD62.12}
\]
both satisfy the monodromy condition, and their special characters
on prime-to-$p$ inertia differ.  The group $M$ individually acquires
good supersingular reduction; the stated filtration supplies the
monodromy condition for $H$.
Consequently a universal B.4 index in height $h\ge2$ satisfies
$\nu_{\min}(p,e,h)\ge n$ in this range.
\end{theorem}
\begin{proof}
Take $G$ and its rational point $P$ from the arc lemma, and put
$C=\langle P\rangle\simeq\mathbf Z/p^n$ and $G'=G/C$.
The finite-flat closure of $C$ in the good model over $\mathcal O_L$
and its quotient show that $G'$ also acquires good reduction.
Equivalently this is the usual quotient of a BT group by a finite
flat subgroup; the quotient is isogenous to the original group.

The rank-one crystalline determinant argument in the full proof of
Proposition~\ref{prop:B-RD18}, using
\cite[Corollary~2.2.6]{Kisin2006Crystalline}, gives
\[
                      \det T_pG=\mathbf Z_p(1)\otimes\psi,
\]
where $\psi$ is trivial on $G_L$.  Indeed, over $L$ the determinant
has weight one and is the Tate representation of a height-one
multiplicative group; the residue field is algebraically closed,
so there is no residual unramified character.  Its finite inertia
value on $\sigma$ is
\[
                         \psi(\sigma)=a^{p+1}.
                                                        \tag{B-RD62.13}
\]
Thus $\psi$ is a character of order $p-1$, valued in the Teichmuller
units of $\mathbf Z_p$.

The multiplication-by-$p^n$ snake sequence for $G\to G'$ is
\[
 0\longrightarrow C\longrightarrow G[p^n]
  \longrightarrow G'[p^n]\longrightarrow C\longrightarrow0.
\]
The image is the cyclic quotient $G[p^n]/C$.  Because $C$ is
constant and primitive as a free $\mathbf Z/p^n$ submodule, the
determinant identifies this quotient with
$(\mathbf Z/p^n)(\chi_p\psi)$.  Set $M=G'\otimes\psi^{-1}$.
It follows that its actual truncation fits into
\[
 0\longrightarrow\mu_{p^n}\longrightarrow M[p^n]
       \longrightarrow(\mathbf Z/p^n)(\psi^{-1})
       \longrightarrow0.
\]
Lemma~\ref{lem:B-RD62-extension} constructs $H$ with precisely this
truncation.  Twisting by the finite character preserves potential
good reduction of $M$; $H$ has the explicitly stated good graded
pieces.

Isogeny preserves the special rational character, and the finite
twist acts on it with the same scalar.  Thus, at the chosen generator
of the finite tame descent quotient, the two eigenvalue multisets are
\[
 \begin{array}{c|c}
 M&\{a^{-p},a^{-1}\}=\{\zeta^p,\zeta\}\\
 H&\{1,a^{-(p+1)}\}=\{1,\zeta^{p+1}\}.
 \end{array}                                           \tag{B-RD62.14}
\]
Their trace difference is
\[
 \zeta+\zeta^p-1-\zeta^{p+1}
                  =-(\zeta-1)(\zeta^p-1)\ne0.
\]
The generator has order $N$ in the finite special-inertia quotient,
which is prime to $p$; no finite-order lift in absolute inertia is
asserted.  Adding the same constant etale group of height $h-2$
preserves the truncation isomorphism and the trace difference.
Since $p^n$ is the manuscript level $\nu=n-1$, failure there gives
the lower bound $\nu_{\min}\ge n$.
\end{proof}

\begin{corollary}[The first ramification step and an infinite sequence]
\label{cor:B-RD62}
\status{Proved; exact lower bounds only, with full independent audits.}
For every $p\ge5$ and $h\ge2$,
\[
                      \nu_{\min}(p,2(p-1),h)\ge2.
                                                        \tag{B-RD62.15}
\]
More generally, for $n\ge2$ put $e_n=p^{2n-3}+p-2$.  Then
$p-1\mid e_n$, $p\nmid e_n$, and
\[
                         \nu_{\min}(p,e_n,h)\ge n.
                                                        \tag{B-RD62.16}
\]
\end{corollary}
\begin{proof}
The integer $e_n$ is the least multiple of $p-1$ strictly greater
than $p^{2n-3}$.  It is strictly smaller than $p^{2n-2}$, and its
residue modulo $p$ is $p-2\ne0$.  Apply the theorem.  For $n=2$
it is $2(p-1)$.
\end{proof}

\begin{remark}[Scope and exact finite checks]
\label{rem:B-RD62-scope}
\status{Proved scope and limitation; no upper bound or exact next index.}
These examples compare a potentially good group with a group having
a two-step potentially good filtration.  They do not assert that
both entire groups are potentially good.  They are compatible with
Theorem~\ref{thm:B-RD61}: the smallest new ramification is $2(p-1)$,
whereas that theorem has $e=p-1$.  The formal group and its exact
descent arc are specified by (B-RD62.4)--(B-RD62.6); the filtered
comparison uses the explicit cohomological lifting argument.
No finite computation is being used to assert effectivity of either
construction.  The exact certificate
\path{code/verify_tame_cyclic_level_counterexamples.py}
checks $15{,}556$ arithmetic cases, $94{,}756$ strict iterate-weight
comparisons, $14{,}592$ logarithm monomials, $40$ independently
recursively computed multiplication coefficients, $24$ finite
character identities, and $264$ first-threshold calculations.
The range is primes $5\le p\le101$ and levels $2\le n\le12$:
all admissible $e$ are tested for $n=2$, and explicitly specified
first, last, and middle windows are tested at higher levels.
The direct multiplication coefficients through degree $p^2$ are
checked for $p=5,7,11,13$.  The actual author run on
19 November 2026 took $0.129167474$ seconds and passed; the complete
source and output are retained.  Full independent root and obstruction/crystal
audits passed, with author-preserving exact reruns taking
$0.137289279$ and $0.123158989$ seconds, respectively.
The upper bound at $e=2(p-1)$
and the general least-index problem remain open here.
\end{remark}
